% Generated draft: requires mathematical and source review.
\documentclass[11pt]{article}
\usepackage[margin=25mm]{geometry}
\usepackage{fontspec}
\setmainfont{Latin Modern Roman}
\usepackage{amsmath,amssymb,mathtools,unicode-math}
\setmathfont{Latin Modern Math}
\usepackage{enumitem,needspace,xurl,hyperref}
\setlength{\parindent}{0pt}
\setlength{\parskip}{5pt}
\setlength{\emergencystretch}{3em}
\newcommand{\N}{\mathbb{N}}
\newcommand{\Z}{\mathbb{Z}}
\newcommand{\Q}{\mathbb{Q}}
\newcommand{\R}{\mathbb{R}}
\newcommand{\C}{\mathbb{C}}
\newcommand{\F}{\mathbb{F}}
\newcommand{\A}{\mathbb{A}}
\newcommand{\PP}{\mathbb{P}}
\newcommand{\E}{\mathbb{E}}
\newcommand{\eps}{\varepsilon}
\newcommand{\dd}{\,\mathrm{d}}
\begin{document}
\section*{UnsolvedMath ID: 700192}
\textbf{Source problem code:} catalogue entry
\subsection{Definitions and mathematical statement}
Let \(\mu\) be an infinite cardinal, identified with its initial ordinal. To make the source's terminology explicit, adopt the usual continuous transfinite-series convention: a group \(H\) is \(\mu\)-overnilpotent if, for every \(h\in H\), there exist an ordinal \(\lambda<\mu\) and subgroups \((H_\alpha)_{\alpha\leq\lambda}\) satisfying
\[
H_0=\langle h\rangle,\qquad H_\lambda=H,\qquad
H_\alpha\trianglelefteq H_{\alpha+1}\quad(\alpha<\lambda),
\]
with the family ascending and
\[
H_\delta=\bigcup_{\alpha<\delta}H_\alpha
\quad\text{for every nonzero limit ordinal }\delta\leq\lambda.
\]
Thus each cyclic subgroup begins a possibly transfinite normal series reaching \(H\). Normality here means normality in the next subgroup. The source does not explicitly specify continuity or its convention for measuring series length; these conventions require confirmation.

Write \(\mathcal R_\mu\) for the class of \(\mu\)-overnilpotent groups. The source states that \(\mathcal R_\mu\) is a radical class. Under the standard radical-class interpretation, call \(G\) \(\mu\)-semisimple if
\[
\forall N\trianglelefteq G,\qquad
N\in\mathcal R_\mu\ \Longrightarrow\ N=\{1\}.
\]
This interpretation of semisimplicity is not explicitly defined in the supplied record.

The question attributed to B. I. Plotkin (1969) is whether
\[
\forall\text{ infinite cardinals }\mu_1<\mu_2\quad
\exists\text{ a group }G:\qquad
G\in\mathcal R_{\mu_2}
\quad\text{and}\quad
\forall N\trianglelefteq G\,
\bigl(N\in\mathcal R_{\mu_1}\Longrightarrow N=\{1\}\bigr).
\]
The supplied wording does not explicitly require \(G\neq\{1\}\). Whether nontriviality is intended must be resolved before treating this as a precise separation problem; it has not been silently added to the displayed question.

The source's \emph{Remark of 2001} attributes to S. Vovsi (1972) the weaker assertion
\[
\forall\text{ infinite cardinals }\mu_1<\mu_2\quad
\exists G\in\mathcal R_{\mu_2}\setminus\mathcal R_{\mu_1}.
\]
This supplied historical assertion is distinct from the requested semisimplicity condition. The record also reports an affirmative answer as an explicitly unverified AI solution claim; it does not supply a verified resolution.
\subsection{Short English statement}
For every pair of infinite cardinals \(\mu_1<\mu_2\), Plotkin asks for a \(\mu_2\)-overnilpotent group having no nontrivial normal \(\mu_1\)-overnilpotent subgroup. The intended nontriviality of the group and the precise transfinite-series conventions need confirmation. The supplied 2001 remark reports a weaker result, and the reported 2026 affirmative claim is explicitly unverified.
\subsection{Sources}
\begin{itemize}
\item Local UnsolvedMath registry: permanent ID \(700192\), exact source problem code \texttt{catalogue entry}, title \emph{Kourovka Notebook Problem 3.43}; supplied complete source record, upstream revision \texttt{d371ee917ac9647fbb06cd6fc529deec6f3dda82}.
\item B. I. Plotkin, Problem 3.43, \emph{The Kourovka Notebook (Unsolved Problems in Group Theory)}, 3rd issue (1969). The supplied record identifies the collection as Issues 1--21, edited by E. I. Khukhro and V. D. Mazurov, and cites the 21st issue edition (2026), \texttt{arXiv:1401.0300v48}.
\item S. Vovsi, \emph{Sov. Math. Dokl.} \textbf{13} (1972), 408--410. The supplied \emph{Remark of 2001} attributes the weaker strict-inclusion result to this reference; no article title is supplied.
\item Achyuth Jayadevan, \emph{Preprint} (2026), \url{https://zenodo.org/records/22721386}. The supplied background and status-review notes describe this as an unverified AI solution claiming an affirmative answer, marked with \(\#\) in the October 2026 21st issue; no preprint title is supplied.
\item Supplied status curation dated \texttt{2026-10-16}: the registry status and canonical status are \texttt{open}, with evidence classified as \texttt{unrefereed\_claim}. This is supplied dated curation, not an independently verified literature assessment. None of the cited sources has been independently checked for this draft.
\end{itemize}
\end{document}
